\section{实战复盘一：三类高频任务如何跑通端到端闭环}

为了避免把访谈停留在观点层，这一章把 OpenClaw 的思路落到可执行任务上。目标不是“复刻某个 demo”，而是理解一个 agent 在真实业务中需要跨越哪些环节。

\subsection{任务 A：发票与对账自动化}

访谈末段提到小企业通过 OpenClaw 自动化繁琐流程，其中最典型的是发票收集、分类、回填与提醒。这类任务天然适合 agent，因为它跨多个系统且重复率高。

\begin{knowledgebox}{发票自动化的标准链路}
\begin{enumerate}
    \item 从邮箱或云盘抓取新票据；
    \item 用 OCR + 规则提取金额、税率、日期、供应商；
    \item 与账务系统主数据对齐；
    \item 对异常项触发人工确认；
    \item 生成每周对账摘要并推送负责人。
\end{enumerate}
\end{knowledgebox}

这一链路里，模型并不负责“算账正确性”，而负责\textbf{把异构信息串起来}。真正决定可靠性的，是 schema 约束、字段校验、异常兜底和可回溯日志。

\begin{importantbox}{把“自动化”做成“可审计自动化”}
若系统只输出结果、不保留中间证据，企业无法通过审计。建议每个关键动作记录：
\begin{itemize}
    \item 输入来源（文件哈希、邮件 ID）；
    \item 转换步骤（提取模型版本、规则版本）；
    \item 输出去向（写入 API、表单、通知渠道）；
    \item 人工介入点（谁在何时确认了什么）。
\end{itemize}
\end{importantbox}

\subsection{任务 B：个人助理型任务编排}

访谈中大量讨论“我要一句话让 agent 完成多步任务”，例如行程提醒、联系人通知、服务下单。它的难点不是 API 调用本身，而是跨上下文的一致性。

\begin{warningbox}{个人任务编排最常见失败模式}
\begin{itemize}
    \item 时间上下文错位：时区、日期歧义、重复提醒；
    \item 身份上下文错位：同名联系人误触达；
    \item 意图漂移：从“提醒”变成“自动执行购买”；
    \item 权限升级：低风险请求意外触发高风险动作。
\end{itemize}
\end{warningbox}

解决方法是把任务拆为“计划层 + 执行层”：计划层先产出明确步骤并等待确认，执行层再按预算（time/cost/permission）完成动作。这样可以把“会做”与“该不该做”分离。

\subsection{任务 C：客服与工单响应增强}

对小团队而言，客服响应常常是低价值高重复劳动。OpenClaw 这一类 agent 可以先把工单分类、草拟回复、补充上下文，再交给人类审核发送。

\begin{knowledgebox}{客服场景的最小质量指标}
\begin{itemize}
    \item 首次响应时间（FRT）；
    \item 自动草拟可采纳率；
    \item 二次修改幅度；
    \item 升级到人工专家的触发准确率。
\end{itemize}
如果仅看“回复速度”，很容易得到高效但错误率高的系统。
\end{knowledgebox}

\subsection{为什么这些任务都依赖同一个底层能力}

无论是财务、个人助理还是客服，真正共通的是三件事：稳定上下文、可控权限、可审计执行。模型能力只决定上限，系统治理决定下限。

\begin{importantbox}{端到端闭环的评价公式}
可把 agent 任务质量粗略写成：
\[
Q = \alpha \cdot U + \beta \cdot R + \gamma \cdot S
\]
其中 \(U\) 是用户价值（节省时间、提升体验），\(R\) 是可靠性（正确率、可恢复性），\(S\) 是安全性（权限约束、攻击韧性）。产品化要同时优化三项，而非单点最大化。
\end{importantbox}

\subsection{本章小结}
访谈中的“好用”来自端到端闭环，而不是某个模型参数。能否把任务拆成稳定链路，是 OpenClaw 类系统从 demo 走向业务价值的关键。

